\scanpage{71}
% loc: 1–2
Donc $b^2=1$ i.e. $b=\pm1$ donc
\[
 \sigma=\pm(1-2\rho)
\]

% loc: 3–4
\deleted{Supposons qu'on ait}
\[
 \sigma=1-2\rho
\]

% loc: 5–8
alors on ne peut avoir en $=$ temps $-(1-2\rho)$ (car $2\,1_k\ne0$), donc cela fixe $\rho$ dans la \deleted{paire} $\{\rho,\rho'=1-\rho\}$. Mais il faut voir lequel ! \deleted{Donc quitte à rectifier le tir, écrivons}

% loc: 9, crossed formula
\[
 \deleted{u+v+w=1-2\rho\qquad \rho\in M^*,\ \det(\rho)=1,\ \rho^3=1}
\]

% loc: 10–15
Revenant : $M$, on peut dire que pour un triangle $\{u,v,w\}$ de $C^*(e)$, existe $\rho\in M^*$ avec
\[
 \left\{\begin{aligned}
 \det\rho&=1&\rho&\ \text{induit une permutation}\\
 &&&\text{circulaire entre les }u,v,w\\
 \rho^3&=1&&\text{et on aura }\operatorname{Tr}\rho=\operatorname{Tr}\rho'=1
 \end{aligned}\right.
\]

% loc: left of conditions, oblique
% slanted
\margin{C'est la relation $\operatorname{Tr}\rho=1$ qui est la plus simple pour lever l'ambiguité mod scalaire (une fois que la perm. \unsure{ass.} est fixée), \deleted{\unsure{dét}} \unsure{les rel.} $\det\rho=1$, $\rho^3=1$ viennent par surcroît $\ldots$}

% loc: 16–18
\unsure{que|qui} $\rho$, à passage à $\rho'$ près. Il y a cependant un seul tel $\rho$, tel que
\[
 \boxed{u+v+w=1-2\rho}
\]
\margin{vérifier si c'est le bon !}

% loc: 19–22, bottom
D'où la structure \deleted{\unsure{de}} orientée : \inserted{pour} $\{u,v,w\}$ formant triangle, \deleted{\unsure{le}} le $\rho$ qu'on vient de décrire définit sur l'ens $\{u,v,w\}$ une permutation circulaire canonique.
% archive sheet number: 70

